\chapter{Lists, and thinking recursively}
\label{ch:lists}

Lisp stands for \emph{list processing}, and Scheme is a Lisp. Nearly every
value \nodwyn{} hands your code is a list: the model's reply, a batch of tool
calls, a tool's result, a session's events. Learn lists well and the rest of
the book reads easily.

\section{Pairs all the way down}

A \textbf{pair} glues two values together. \code{cons} builds one,
\code{car} takes the first half and \code{cdr} the second (the names are
from a 1950s computer; everyone just says ``car'' and ``could-er''):

\begin{scm}
(cons 1 2)                  ;=> (1 . 2)
(car (cons 1 2))            ;=> 1
(cdr (cons 1 2))            ;=> 2
\end{scm}

A \textbf{list} is a chain of pairs whose last \code{cdr} is the empty
list, \code{'()}:

\begin{txt}
   (basil mint sage)

   [basil| ]--->[mint| ]--->[sage| ]---> ()
\end{txt}

\begin{scm}
(cons 1 '(2 3))             ;=> (1 2 3)
(car '(basil mint sage))    ;=> basil
(cdr '(basil mint sage))    ;=> (mint sage)
(cadr '(basil mint sage))   ;=> mint        the car of the cdr
(null? '())                 ;=> #t
(pair? '(1 . 2))            ;=> #t
(list? '(1 . 2))            ;=> #f          no () at the end
\end{scm}

Lists in \nodwyn{} Scheme never change once made: there is no
\code{set-car!}. You build new lists instead, which turns out to be exactly
what you want when a strategy wraps another strategy: nobody can pull the
rug from under anybody.

\section{The everyday toolkit}

\begin{scm}
(list 1 (+ 1 1) 3)                   ;=> (1 2 3)
(append '(1 2) '(3) '() '(4 5))      ;=> (1 2 3 4 5)
(reverse '(1 2 3))                   ;=> (3 2 1)
(length '(a b c))                    ;=> 3
(list-ref '(a b c) 2)                ;=> c          counting from 0
(list-tail '(a b c d) 2)             ;=> (c d)
(iota 5)                             ;=> (0 1 2 3 4)
(member "mint" '("basil" "mint" "sage"))   ;=> ("mint" "sage")
(memq 'x '(a b))                     ;=> #f
\end{scm}

\section{Functions that eat lists}

These four will carry most of your code:

\begin{scm}
(map (lambda (p) (string-append p "!")) '("basil" "mint"))
                                     ;=> ("basil!" "mint!")
(map + '(1 2 3) '(10 20 30))         ;=> (11 22 33)
(filter odd? (iota 10))              ;=> (1 3 5 7 9)
(fold + 0 '(1 2 3 4))                ;=> 10
(sort '("mint" "basil" "sage") string<?)
                                     ;=> ("basil" "mint" "sage")
\end{scm}

\code{fold} walks a list left to right carrying an \emph{accumulator}:
\code{(fold f init list)} calls \code{(f acc x)} for each element. Folding
with \code{cons} reverses a list, which is a nice way to see it working:

\begin{scm}
(fold (lambda (acc x) (cons x acc)) '() '(1 2 3))   ;=> (3 2 1)
\end{scm}

\code{for-each} is \code{map} for effects: it calls the function on each
element and throws the results away. \code{apply} calls a function with a
list as its arguments:

\begin{scm}
(apply max '(3 9 4))                 ;=> 9
(apply + 1 2 '(3 4))                 ;=> 10    loose arguments, then a list
\end{scm}

\section{Thinking recursively}

To handle a list, handle the empty list, and handle ``the first element and
the rest''. The rest is a smaller list, so the same function handles it:

\begin{scm}
(define (sum l)
  (if (null? l)
      0
      (+ (car l) (sum (cdr l)))))
(sum '(1 2 3 4))                     ;=> 10
\end{scm}

That is how \code{map} itself could be written:

\begin{scm}
(define (my-map f l)
  (if (null? l)
      '()
      (cons (f (car l)) (my-map f (cdr l)))))
(my-map (lambda (x) (* 10 x)) '(1 2 3))    ;=> (10 20 30)
\end{scm}

And recursion handles shapes that loops find awkward, like lists inside
lists:

\begin{scm}
(define (flatten x)
  (cond ((null? x) '())
        ((pair? x) (append (flatten (car x)) (flatten (cdr x))))
        (else (list x))))
(flatten '(1 (2 (3 4)) ((5))))       ;=> (1 2 3 4 5)
\end{scm}

\section{Tail calls: loops for free}

Look at \code{sum} again. After the recursive call returns, there is still
work to do: add \code{(car l)}. Each call waits for the next, so a list a
million long needs a million waiting calls.

Now look at this:

\begin{scm}
(define (count-up n)
  (let loop ((i 0) (acc 0))
    (if (> i n)
        acc
        (loop (+ i 1) (+ acc i)))))
(count-up 1000000)                   ;=> 500000500000
\end{scm}

The recursive call to \code{loop} is the \emph{last thing} the function
does: it is in \textbf{tail position}. Nothing waits for it, so Scheme
reuses the current call instead of stacking a new one. This is a
\textbf{proper tail call}, and it is guaranteed: a loop written this way
runs in constant space, whether it goes around ten times or ten billion.
The last expression of a \code{lambda}, a \code{begin}, a \code{let}
body, a \code{cond} clause, \code{and}/\code{or} and both branches of
\code{if} are all tail positions.

\code{(let loop ((i 0) (acc 0)) \ldots)} is a \textbf{named let}. It
defines a local function \code{loop} with those parameters and calls it
straight away with those starting values. It is the idiom for loops in
Scheme, and you will see it on almost every page of \nodwyn's own
\file{internal/loop/*.scm}.

\begin{careful}
Non-tail recursion is fine for trees and short lists, but it uses the Go
stack, and a recursion several hundred thousand calls deep exhausts it
and ends the whole process rather than raising an error you can catch. When
a list could be long (a log, a file's lines, a directory listing), write the
loop in tail position with an accumulator, as \code{count-up} does.
\end{careful}

\section{Association lists: little dictionaries}

A list of pairs is the simplest dictionary there is, and \nodwyn{} uses it
everywhere: tool calls, results, a model's reply, tool metadata, parsed
JSON.

\begin{scm}
(define plants '(("basil" . "daily") ("mint" . "weekly")))
(assoc "mint" plants)                ;=> ("mint" . "weekly")
(assoc-ref plants "mint")            ;=> "weekly"
(assoc-ref plants "rose" "unknown")  ;=> "unknown"
(assq 'b '((a . 1) (b . 2)))         ;=> (b . 2)
\end{scm}

\code{assoc-ref} is the one you will use most. It takes the list, the key
and an optional default. \code{assoc} compares keys with \code{equal?}
(good for strings), \code{assq} with \code{eq?} (fine for symbols).

\begin{tip}
A model's reply comes to your strategy as an association list with
\emph{symbol} keys: \code{(assoc-ref reply 'tool-calls '())}. A parsed JSON
object has \emph{string} keys: \code{(assoc-ref args "path")}. Mixing them
up is the most common bug in a first strategy, and the symptom is a
default value where you expected data.
\end{tip}

\section{Loops with \texttt{do}}

For the occasional counting loop, \code{do} names its variables, their
starting values and how each steps, then the test that ends it and what to
return:

\begin{scm}
(do ((i 0 (+ i 1))
     (acc '() (cons i acc)))
    ((= i 3) acc))                   ;=> (2 1 0)
\end{scm}

Anything after the test clause runs once per round, before the step.
Named \code{let} is usually clearer; \code{do} shines when you fill a
vector (Chapter~\ref{ch:playground}).

\begin{puzzle}[ 3: the harvest]
Given \code{(define harvest '(("tomato" . 12) ("basil" . 3) ("bean" .
40)))}, write \code{(total harvest)} that adds up the counts with
\code{fold}, and \code{(biggest harvest)} that returns the name with the
largest count. Then write \code{total} again with a named \code{let}.\ans{3}
\end{puzzle}

\begin{puzzle}[ 4: your own filter]
Write \code{(keep pred l)} that behaves like \code{filter}, first with plain
recursion, then in tail position with an accumulator (you will need
\code{reverse} at the end). Why the \code{reverse}?\ans{4}
\end{puzzle}

\begin{puzzle}[ 5: zip]
Write \code{(zip a b)} so that \code{(zip '(1 2 3) '(x y z))} is
\code{((1 . x) (2 . y) (3 . z))}. There is a one-line answer using
\code{map}.\ans{5}
\end{puzzle}
